\SPBGContentSlide
  {SMO против библиотечных SVM}
  {%
    \centering
    \TinyBody
    \renewcommand{\arraystretch}{1.12}
    \setlength{\tabcolsep}{4pt}
    \begin{tabular}{|l|l|r|c|l|}
      \hline
      \textbf{Блок} & \textbf{Метод} & \textbf{$\Delta F_1$ к SMO} & \textbf{$p_H$} & \textbf{Вывод} \\
      \hline
      N1--N4 & SVC & 0.000 & 1.000 & качество совпадает \\
      \hline
      N1--N4 & LinearSVC & -0.057 & 0.004 & ниже SMO \\
      \hline
      R1--R15 & SVC & +0.028 & 0.036 & выше SMO \\
      \hline
      R1--R15 & LinearSVC & +0.031 & 0.036 & выше SMO \\
      \hline
    \end{tabular}

    \vspace{.16cm}
    \begin{tabular}{|l|l|r|c|}
      \hline
      \textbf{Блок} & \textbf{Метод} & \textbf{$\Delta t$ к SMO, c} & \textbf{$p_H$} \\
      \hline
      N1--N4 & SVC & -0.0344 & $3.51\cdot10^{-7}$ \\
      \hline
      N1--N4 & LinearSVC & -0.0358 & $3.51\cdot10^{-7}$ \\
      \hline
      R1--R15 & SVC & -0.0369 & $1.42\cdot10^{-12}$ \\
      \hline
      R1--R15 & LinearSVC & -0.2330 & $5.48\cdot10^{-16}$ \\
      \hline
    \end{tabular}

    \vspace{.16cm}
    \begin{minipage}{.90\linewidth}
      \SmallBody
      \raggedright
      На синтетических неразделимых задачах собственный SMO воспроизводит
      качество SVC, но на реальных данных уступает библиотечным решателям.
      Главная причина видна на крупных и несбалансированных наборах: требуется
      более сильная эвристика выбора рабочей пары и контроль лимита итераций.
    \end{minipage}
  }
